\textbf{The two-tier saturation picture.}
\textbf{(a)} The V-axis is encoded at two scales in a converged EEG
classifier. \emph{Tier 1} is the shared class-mean basin: the nine
emotion centroids project onto a common one-dimensional valence axis,
and this projection is consistent across checkpoints (the
``saturated'' tier---adding a V-axis loss cannot improve a structure
the network already encodes from task labels). \emph{Tier 2} is the
within-class residual: the per-trial structure that varies between
seeds and is averaged out by ensembling. The $r{=}{+}0.74$
correlation between within-class residual encoding and per-checkpoint
ensemble contribution (Figure~\ref{fig:two-tier-scatter}) lives in
this tier. \textbf{(b)} The saturation transition. Below the
threshold (e.g.\ CBraMod or EMOD-$d{=}3$ vanilla, BACC ${\le}0.62$),
the network has not yet converged to the V-axis basin and adding a
V-axis loss is approximately a no-op. Above the threshold (the strong
$d{=}6+\text{aug}+\text{KD}$ recipe, BACC ${\ge}0.66$), the network
is already in the basin and the same loss perturbs the saturated
component while leaving the load-bearing residual untouched, so
accuracy drops. The $25$ interventions of
Figure~\ref{fig:forest} form the empirical content of the cliff on
the right of \textbf{(b)}.
